% Actual class: 12, Week 12, 2026-11-06.
% Source: 12_gan.pdf pp. 1--91; five complete English caption tracks.
% Prepared from templates/lecture.tex; no full video-frame review.

\section{第 12 节课：生成对抗网络}
\label{lecture:SC4001-L12}

\lecturemeta
  {SC4001-L12}
  {2026-11-06（第十二周）}
  {12\_gan：Generative Adversarial Networks}
  {pp.~1--91；两个实验例题 pp.~56--62}
  {12\_gan\_part1--5，约 72.5 分钟；完整英文字幕与讲义核对}
  {已讲解；公式与来源已核对，笔记待审阅}

本讲以生成手写数字为例：输入随机向量，输出新的图片；训练一个判别器提供真假反馈，再用它训练生成器。重点是损失的优化方向、交替更新时的梯度路径，以及生成结果的真实性与多样性。

\lecturetopic{SC4001-L12-T01}{生成任务与两个网络的分工}

\textbf{来源：}pp.~20--27，视频 Part 2；pp.~2--19、Part 1 展示生成、修复、图像转换与编辑等应用。

Generator（生成器）$G_\theta$ 把简单分布中的随机向量映射为数据：
\[
 z\sim p(z),\qquad \hat x=G_\theta(z).
\]
$p(z)$ 常取 Gaussian 或 uniform distribution（高斯或均匀分布），$\theta$ 为生成器参数。不同 $z$ 的输出共同形成生成分布 $p_g$，目标是使 $p_g\approx p_{\rm data}$。训练集提供真实样本，但每个 $z$ 没有指定的配对目标图片。

Discriminator（判别器）$D_\phi(x)\in(0,1)$ 输出将输入判为真实数据的分数，$\phi$ 为其参数。训练 $D$ 时真实样本的标签为 $1$，生成样本的标签为 $0$；训练 $G$ 时，希望 $D(G(z))\to1$。真假标签由数据来源自动确定，vanilla GAN（普通 GAN）不需要类别标签。

\begin{center}
% Original schematic based on 12_gan.pdf pp. 22, 29--30.
\begin{tikzpicture}[
 font=\small,>=Stealth,
 box/.style={draw=noteBlue,fill=noteBlue!4,rounded corners,
   text width=1.8cm,minimum height=0.95cm,align=center,inner sep=4pt},
 flow/.style={->,thick,noteBlue}
]
\node[box] (z) at (0,0) {随机输入\\$z\sim p(z)$};
\node[box] (g) at (3.1,0) {生成器\\$G_\theta$};
\node[box] (fake) at (6.2,0) {生成样本\\$\hat x=G(z)$};
\node[box] (d) at (9.3,0) {判别器\\$D_\phi$};
\node[box] (out) at (12.4,0) {真假分数\\$(0,1)$};
\node[box] (real) at (6.2,1.65) {真实样本\\$x\sim p_{\rm data}$};
\draw[flow] (z)--(g);
\draw[flow] (g)--(fake);
\draw[flow] (fake)--(d);
\draw[flow] (d)--(out);
\draw[flow] (real.east)-|(d.north);
\draw[flow,dashed] (d.south)--++(0,-0.8)-|node[pos=0.5,below,align=center]
  {更新 $G$ 时：梯度经过 $D$，固定 $D$ 的参数}(g.south);
\end{tikzpicture}

\end{center}
{\small 原创示意，依据 pp.~22、29--30；实线为前向计算，虚线为更新生成器时的反向梯度。}

$D(x)$ 是真假分类输出，不是数据密度 $p_{\rm data}(x)$。在真假采样比例相同、$p_g=p_{\rm data}$ 且 $D$ 达到最优时，$D(x)=1/2$；单独观察到 $0.5$ 也可能是判别器尚未学会区分，仍需检查生成分布。

训练完成后只需从 $p(z)$ 抽样并运行 $G$，无需输入真实图片或运行 $D$。普通 GAN 属于 implicit generative model（隐式生成模型）：可以采样，通常不直接给出可计算的 $p_g(x)$。

\lecturetopic{SC4001-L12-T02}{Minimax、交替更新与非饱和损失}

\textbf{来源：}pp.~27--55，视频 Part 3；以下 $\log$ 均指自然对数。

\subsubsection{先明确每个目标的优化方向}

原始 GAN 是 two-player minimax game（双人极小极大博弈）：
\[
 \min_\theta\max_\phi V(D_\phi,G_\theta),\qquad
 V(D,G)=\mathbb E_{x\sim p_{\rm data}}[\log D(x)]
       +\mathbb E_{z\sim p(z)}[\log(1-D(G(z)))].
\]
期望在训练中用批量平均估计。固定 $G$ 时，$D$ 最大化 $V$：提高真实样本分数、降低生成样本分数。若使用最小化优化器，令
\[
 L_D=-\mathbb E_x[\log D(x)]-\mathbb E_z[\log(1-D(G(z)))].
\]
固定 $D$ 时，真实样本项与 $\theta$ 无关，原始生成器损失为
\[
 L_G^{\rm MM}=\mathbb E_z[\log(1-D(G(z)))].
\]
最小化它会提高 $D(G(z))$。讲义的 $0.8/0.2$ 数值比较见\relatedexercise{SC4001-L12-E01}{损失值与优化方向}。

\subsubsection{冻结参数与切断梯度是两件事}

\begin{center}
\small
\begin{tabularx}{\textwidth}{@{}lXX@{}}
\toprule
步骤 & 更新的参数与目标 & 梯度处理 \\
\midrule
训练 $D$ & 只更新 $\phi$；真实为 $1$、生成为 $0$ & 把 $G(z)$ 视为固定输入，可使用 \texttt{detach()} \\
训练 $G$ & 只更新 $\theta$；希望生成样本被判为 $1$ & 固定 $D$ 的参数，保留经过 $D$ 到 $G$ 的梯度 \\
\bottomrule
\end{tabularx}
\end{center}

\textbf{补充梯度表达。}对一个样本，令 $\hat x=G_\theta(z)$、$s=D_\phi(\hat x)$，链式法则给出
\[
 \nabla_\theta L_G
 =\frac{\partial L_G}{\partial s}
   J_{G_\theta}(z)^\top\nabla_{\hat x}D_\phi(\hat x),
\]
其中 $J_{G_\theta}(z)=\partial G_\theta(z)/\partial\theta$。生成器利用判别器对输入的导数调整输出；训练 $G$ 时若对 $G(z)$ 做 \texttt{detach()}，或禁用整个 $D$ 前向过程的梯度记录，就会切断这条路径。训练通常交替执行若干次 $D$ 更新与一次 $G$ 更新，次数比例属于训练设置。

\subsubsection{Non-saturating loss：改善训练初期的梯度}

训练初期生成图片较差，$s=D(G(z))$ 接近 $0$。常用的 non-saturating loss（非饱和损失）为
\[
 L_G^{\rm NS}=-\mathbb E_z[\log D(G(z))].
\]
它同样推动 $s\to1$，但梯度与原始损失不同。设 $s=\sigma(a)$，$a$ 为判别器最后一个 sigmoid 之前的 logit，$\mathrm ds/\mathrm da=s(1-s)$。对单个样本，补全 p.~55 的链式求导：
\[
 \frac{\partial\log(1-s)}{\partial a}
 =-\frac1{1-s}s(1-s)=-s,
 \qquad
 \frac{\partial(-\log s)}{\partial a}
 =-\frac1s s(1-s)=s-1.
\]
当 $s=0.01$ 时，两者分别为 $-0.01$ 与 $-0.99$，因此非饱和目标在这个位置提供更强的梯度。比较对象是\emph{对 logit 的导数}；传到生成器参数的梯度还取决于中间网络的导数。

\lecturetopic{SC4001-L12-T03}{图像生成结构与潜在空间}

\textbf{来源：}pp.~56--73，视频 Part 3--4。

\subsubsection{从全连接网络到 DCGAN}

讲义先用一维噪声生成 Gaussian 数据，再用 $100$ 维噪声生成 $28\times28$ 的 MNIST 图像。生成器输出 $784$ 个像素值，判别器把图片展平为 $784$ 维输入，最后只输出一个真假分数。完整结构见\relatedexercise{SC4001-L12-E02}{一维分布学习}与\relatedexercise{SC4001-L12-E03}{MNIST 生成实验}。

Deep Convolutional GAN（DCGAN，深度卷积生成对抗网络）利用局部空间关系和卷积权重共享构建图像。讲义生成器先把 $z\in\mathbb R^{100}$ \emph{投影后重排}为小特征图，再逐层扩大空间尺寸：
\[
 \begin{aligned}
 z&\xrightarrow{\text{可学习投影与重排}}4\times4\times1024
    \to8\times8\times512\to16\times16\times256\\
  &\to32\times32\times128\to64\times64\times3.
 \end{aligned}
\]
投影会产生新的特征值；仅 reshape 不能把 $100$ 个数变成 $4\times4\times1024$ 个数。判别器沿相反的空间尺度方向提取特征，再输出真假分数。

\subsubsection{转置卷积与 Batch Normalization}

Transposed convolution（转置卷积）是一种可学习的上采样。讲义用“插入零、按需补边，再应用卷积核”解释其计算。p.~67 的中间网格为 $7\times7$，以步幅 $1$、$3\times3$ 核做不补边卷积，输出边长为 $7-3+1=5$。

这里的“转置”指卷积线性算子的转置，不是矩阵求逆；它通过训练学习高分辨率输出，不能自动恢复卷积丢失的信息。DCGAN 在多数中间层使用 batch normalization（批归一化），对每个通道结合批内与空间位置统计量归一化；通常避开生成器输出层和判别器输入层。

\subsubsection{插值与语义方向}

在两个潜在向量间插值：
\[
 z(t)=(1-t)z_1+tz_2,\quad t\in[0,1],
 \qquad \hat x(t)=G(z(t)).
\]
观察输出是否平滑变化，可检查沿这条路径的生成行为。平滑插值只反映所选路径，判断训练样本记忆还需要其他证据。

讲义还展示潜在空间中的近似向量运算：
\[
 d_{\rm glasses}=z_{\rm man+glasses}-z_{\rm man},\qquad
 z_{\rm new}=z_{\rm woman}+d_{\rm glasses}.
\]
$G(z_{\rm new})$ 可表现出戴眼镜的变化。这里相加的是潜在向量；该方向对其他样本是否有效、是否同时改变其他属性，需要分别检查。

\lecturetopic{SC4001-L12-T04}{训练稳定性与模式崩塌}

\textbf{来源：}pp.~74--81，视频 Part 5。

\subsubsection{交替优化可能振荡}

生成器改变假样本分布，判别器改变生成器的训练反馈，因此双方的目标随对方更新而改变。原始 minimax 希望达到 Nash equilibrium（纳什均衡）：固定另一方时，任何一方都无法单独改善自己的目标。交替梯度更新可能振荡或不收敛，结果也受学习率和双方能力平衡影响。

讲义用 $\min_x\max_y xy$ 展示这种相互影响；它是局部动力学示意，不能当作收敛保证。方向分析与补充的离散更新计算见\relatedexercise{SC4001-L12-E04}{双人目标的旋转趋势}。GAN 的 loss 波动需要结合样本与分布判断，单条曲线不足以验收生成质量。

\subsubsection{真实感与多样性分别检查}

Mode collapse（模式崩塌）指生成分布只覆盖真实分布的一部分。例如，不同 $z$ 都生成几乎一样的数字 $7$：单张图可以很逼真，但其他类别和笔迹缺失。实践中，生成器可能反复利用少数能骗过当前判别器的输出。

Mini-batch discrimination（小批量判别）把批内样本之间的相似性特征加入判别器：先提取每个样本的特征，再比较样本关系，让 $D$ 能识别一批生成结果是否过于相似，从而为 $G$ 提供多样性反馈。这是额外的跨样本比较机制，普通小批量训练或 BatchNorm 本身不等同于该机制。

\lecturetopic{SC4001-L12-T05}{条件控制：cGAN 与 AC-GAN}

\textbf{来源：}pp.~82--85，视频 Part 5；pp.~86--89 为本讲总结，pp.~90--91 为参考资料。

Conditional GAN（cGAN，条件生成对抗网络）学习 $p_{\rm data}(x\mid c)$。条件同时提供给 $G$ 和 $D$：前者生成指定内容，后者判断样本是否符合该条件下的真实分布。

\begin{center}
\small
\begin{tabularx}{\textwidth}{@{}lXX@{}}
\toprule
模型 & 生成器 & 判别器一侧 \\
\midrule
Vanilla GAN & $G(z)$，没有显式类别输入 & $D(x)$ 判断真假 \\
Conditional GAN & $G(z,c)$，给定条件 $c$ & $D(x,c)$ 结合条件判断真假 \\
AC-GAN & $G(z,c)$，给定类别 $c$ & 真假输出 $D(x)$，另加类别输出 $C(x)=P(c\mid x)$ \\
\bottomrule
\end{tabularx}
\end{center}

固定 $c=7$、改变 $z$，可生成不同笔迹的 $7$；类别控制仍需配合类内多样性检查。

Auxiliary Classifier GAN（AC-GAN，辅助分类器 GAN）让判别器同时预测真假和类别。MNIST 的类别输出为十类概率，真假输出仍是二分类。生成器同时接受“看起来真实”和“属于指定类别”两方面的训练信号。

\subsection{一分钟回顾}

GAN 用判别器提供学习信号，使随机输入产生的数据接近真实分布。训练 $D$ 时固定生成样本；训练 $G$ 时固定 $D$ 的参数，但保留经过它的梯度。非饱和损失改善训练初期的反馈，DCGAN 用卷积层次生成图片。生成结果要兼顾真实感与多样性；cGAN 用条件控制生成，AC-GAN 进一步加入类别预测。
